\documentclass{article}
\usepackage[margin=1in]{geometry}
\usepackage{mathtools, amsfonts, amsthm, xcolor, listings}

\title{HW 5}
\author{Sam Ly}

\begin{document}
\maketitle

\section*{Chapter 10}

https://github.com/samlyme/Assignments/tree/cs4080/ch10/challenge/CS4080

\begin{enumerate}
    \item {
        Our interpreter carefully checks that the number of arguments passed to a function matches the number of parameters it expects. Since this check is done at runtime on every call, it has a performance cost. Smalltalk implementations don’t have that problem. Why not?

        Smalltalk doesn't have this problem because you don't actually specify a 
        ``name'' for methods in Smalltalk. Function calling is implemented as a 
        message passing system, so the shape of the message itself determines what
        code is called.
    }

    \item {
        Lox’s function declaration syntax performs two independent operations. It creates a function and also binds it to a name. This improves usability for the common case where you do want to associate a name with the function. But in functional-styled code, you often want to create a function to immediately pass it to some other function or return it. In that case, it doesn’t need a name.

Languages that encourage a functional style usually support anonymous functions or lambdas—an expression syntax that creates a function without binding it to a name. Add anonymous function syntax to Lox so that this works:

        \begin{verbatim}
fun thrice(fn) {
  for (var i = 1; i <= 3; i = i + 1) {
    fn(i);
  }
}

thrice(fun (a) {
  print a;
});
// "1".
// "2".
// "3".
        \end{verbatim}

        First, we create a new Expr type in the GenerateAST class called `Lambda'.
        Then, we create a new class `LoxLambda' which is similar to the existing
        LoxFunction class, just without the `name' field. After that, in the interpreter, 
        visit the lambda expression by returning an instance of `LoxLambda'.
        \begin{verbatim}
    @Override
    public Object visitLambdaExpr(Expr.Lambda expr) {
        return new LoxLambda(expr.params, expr.body, environment);
    }
        \end{verbatim}

        The trickiest part of this challenge was the parser, because of the 
        weird expression-statement case for anonymous functions. As an aside,
        this syntax literally does nothing, even if the params and body are defined,
        the lambda returned is instantly discarded. But the way I approached this 
        was to create a new expression parser function `lambda' that is basically 
        the same as the existing `function' statement parser.
        \begin{verbatim}
    private Expr expression() {
        return lambda();
   }

   private Expr lambda() {
        if (!match(FUN))  return assignment();

       consume(LEFT_PAREN, "Expect '(' after lambda start.");
       List<Token> parameters = new ArrayList<>();
       if (!check(RIGHT_PAREN)) {
           do {
               if (parameters.size() >= 255) {
                   //noinspection ThrowableNotThrown
                   error(peek(), "Can't have more than 255 parameters.");
               }

               parameters.add(
                       consume(IDENTIFIER, "Expect parameter name."));
           } while (match(COMMA));
       }
       consume(RIGHT_PAREN, "Expect ')' after parameters.");

       consume(LEFT_BRACE, "Expect '{' before lambda body.");
       List<Stmt> body = block();
       return new Expr.Lambda(parameters, body);
   }
        \end{verbatim}

        Then, when parsing for declaration, we do this check for the weird edge case
        \begin{verbatim}
    private Stmt declaration() {
        try {
            if (match(FUN)) {
                if (check(IDENTIFIER)) return function("function");
                current--;
            }
            if (match(VAR)) return varDeclaration();
            return statement();
        } catch (ParseError error) {
            synchronize();
            return null;
        }
    }
        \end{verbatim}
        If the name isn't there, then the execution falls through and treats the 
        `fun' keyword as the start of a lambda expression.
    }

    \item {
        Is this program valid?

        \begin{verbatim}
fun scope(a) {
  var a = "local";
}
        \end{verbatim}

        In other words, are a function’s parameters in the same scope as its local variables, or in an outer scope? What does Lox do? What about other languages you are familiar with? What do you think a language should do?

        This is perfectly valid because our implementation treats the parameter 
        names almost like an inherited scope variable. With our block scoping, 
        we can shadow already-defined variables, so this is no different. In my 
        opinion, the way Lox handles this is how languages should handle this anyway.
        Maybe this is bad code style, but that is easily solved by linters in production
        languages.
    }
\end{enumerate}

% \section*{Chapter 11}

\end{document}
